\documentclass{article}%
\usepackage[T1]{fontenc}%
\usepackage[utf8]{inputenc}%
\usepackage{lmodern}%
\usepackage{textcomp}%
\usepackage{lastpage}%
\usepackage{geometry}%
\geometry{margin=2.2cm}%
%
\usepackage{amsmath}%
\usepackage{amssymb}%
\usepackage{booktabs}%
\usepackage{longtable}%
\usepackage{array}%
\usepackage{hyperref}%
\hypersetup{hidelinks}%
\title{ITA 2027 {-} Física (claude\_haiku\_5\_5) x Poliedro\\[4pt]\large Comparação com o gabarito Poliedro Resolve (ITA 2027, 2ª fase, dia 3 {-} Física)\thanks{Gerado com gpt-resolve: \url{https://github.com/lgabs/gpt-resolve}}}%
\date{09/11/2026}%
%
\begin{document}%
\normalsize%
\maketitle%
\section*{Resumo}\textbf{Nota média: 9,5/10} (23 de 24 itens com a resolução correta; cada item vale o mesmo dentro da questão).\par\medskip%
\begin{longtable}{@{}c c p{0.78\textwidth}@{}}\toprule Questão & Nota & Resumo \\ \midrule
1 & 10,0 & A resolução e o gabarito concordam nos três itens. (a) Os dois dão k\_C = MgL/2. (b) As faixas são as mesmas: para k/k\_C $\geq$ 1, só $\theta$ = 0, estável; para 2/$\pi$ < k/k\_C < 1, $\theta$ = 0 instável e $\pm$$\theta$\_e estáveis, com sen$\theta$\_e/$\theta$\_e = k/k\_C; para k/k\_C $\leq$ 2/$\pi$, só $\theta$ = 0, instável. (c) As formas de U($\theta$) são as mesmas; as expressões diferem apenas pela constante Mg$\cdot$h\_O. Conferi numericamente os valores extras da resolução ($\theta$\_e e barreira para $\lambda$ = 0,8, limite da barreira para $\lambda$$\to$1, limiar 8/$\pi$$^2$, coordenadas dos gráficos) e estão todos certos. Os pontos que levantei são de redação e de convenção, e nenhum muda a nota. \\
2 & 10,0 & A resolução e o gabarito do Poliedro coincidem nos dois itens. Em (a), os dois dão H\_n = 0,7\^{}(2n)V$^2$/(8g). Em (b), os dois dão T = 10V/(3g): a energia mínima é a horizontal, (3/8)mV$^2$, atingida no ponto de acumulação dos choques, e o tempo sai da série geométrica dos voos. Conferi o enunciado no PDF: está na página 3 impressa, que é a página 2 do arquivo, e a transcrição está correta. As contas também conferem. Os únicos pontos de dúvida são ambiguidades de redação ('mínimo da sua energia', contagem de n) e erros de digitação na notação do gabarito, nenhum deles capaz de mudar a nota. Questão aprovada nos dois itens. \\
3 & 5,0 & Item a: a resolução e o gabarito dão a mesma distância em formas equivalentes, D = R$\surd$(1 + (2h/g)$\surd$($\mu$e$^2$g$^2$/R$^2$ - $\alpha$$^2$)). Item b: as respostas diferem exatamente por $\surd$(2h/g). A resolução aplica o atrito de Coulomb durante o choque com o solo (impulso $\mu$c$\cdot$m$\cdot$$\surd$(2gh)), e por isso o tempo de queda se cancela: t* = t\_e + $\omega$\_eR/($\mu$c g). O gabarito supõe que a velocidade horizontal se conserva no pouso e chega a t\_e + t\_q + t\_f. Não há erro de conta em nenhum dos dois. Dei razão ao gabarito porque é a leitura pretendida: o enunciado afirma sem condição que a massa desliza, o que não vale sempre no modelo da resolução. O modelo da resolução é fisicamente defensável e merece ser avaliado por um professor. Também tentei registrar essa decisão na memória do time, para que outras correções desta questão sejam consistentes, mas a gravação foi bloqueada pelas restrições de acesso aos arquivos e nada foi salvo. \\
4 & 10,0 & Questão 4: a resolução está correta nos dois itens. No item a, h\_f = h\_i((m+M)/m)\^{}\{3/5\} coincide com o gabarito. No item b, a equação m(h\_max - h\_i) = (3/2)(m+M)h\_i{[}1 - (h\_i/h\_max)\^{}\{2/3\}{]} é equivalente, depois de expandida, à interpretação 1 do Poliedro (adiabática reversível), e a conferência numérica confirma (h\_max $\approx$ 1,455 m no exemplo da resolução). O único ponto aberto é de enunciado: o próprio Poliedro diz que o item b não fixa a evolução do gás. A resolução declara e justifica sua hipótese, que é a leitura mais compatível com 'equação que relaciona h\_max com m, M e h\_i'. \\
5 & 10,0 & A resolução está certa nos dois itens e concorda com o gabarito Poliedro. No item a, as expressões de p' e das posições verticais y' são idênticas às do gabarito, e um traçado de raio exato confirmou ambas. No item b, a resolução dá a fórmula completa $\lambda$0(n1D+n2p)/(2n1n2h). Com p ? D, ela se reduz ao $\lambda$D/(2n2h) do gabarito, e a resolução indica essa aproximação. Um cálculo numérico do caminho óptico exato bate com a fórmula completa, e a do gabarito difere só por um termo de ordem p/D. O único ponto ambíguo é o meio em que $\lambda$ é medido; gabarito e resolução usam o vácuo. Obs.: a questão 5 está na página 6 impressa, que é a página 5 do PDF. Tentei salvar essa nota na memória do time, mas a escrita foi bloqueada nesta sessão. \\
6 & 10,0 & A resolução e o gabarito do Poliedro concordam nos três itens: a) q = 3C$\varepsilon$/4; b) compressão de 2d/3, com as placas a d/3; c) k = 27C$\varepsilon$$^2$/(64d$^2$). Refiz as contas com frações exatas e as três respostas estão certas. A resolução tem um extra: mostra que esse k é exatamente o valor crítico de estabilidade, o que torna o equilíbrio de t3 marginal. Isso é uma sutileza física do enunciado e não altera as respostas. \\
7 & 10,0 & Os dois itens da questão 7 conferem com o gabarito do Poliedro, e a resolução e o gabarito estão corretos. Em (a), |V\_o| = (V\_s/4)|G$\varepsilon$+$\alpha$$\Delta$T|. Em (b), a resposta é a ponte completa com quatro extensômetros: R2 e R3 em tração, R1 e R4 em compressão (ou tudo trocado), com V\_o $\approx$ V\_sG$\varepsilon$ e sem efeito térmico. Conferi o enunciado e a figura no PDF original e refiz as contas numericamente; as formas linear e exata batem. A resolução vai além do gabarito: dá o sinal por braço, testa todas as quantidades de extensômetros (k = 1 a 4) e verifica a forma exata V\_sG$\varepsilon$/(1+$\alpha$$\Delta$T). Não há divergência; ficam só observações de baixa importância sobre a polaridade e a interpretação de 'máxima precisão'. \\
8 & 10,0 & A resolução está certa nos três itens e concorda com o gabarito do Poliedro. a) v = $\surd$(2qV/m), igual ao gabarito. b) Os dois descrevem um movimento helicoidal uniforme: MU com velocidade v cos$\theta$ ao longo de B, composto com MCU de raio mv sen$\theta$/(qB) e período 2$\pi$m/(qB). As justificativas também coincidem. c) Os dois chegam a cos$\theta$ = qBL/(2$\pi$nmv) = (BL/(2$\pi$n))$\surd$(q/(2mV)), com n $\in$ $\mathbb{N}$*. A única diferença está na restrição sobre n. O gabarito escreve cos$\theta$ $\leq$ 1. A resolução escreve 0 < cos$\theta$ < 1, porque o enunciado exige que a partícula saia da linha e volte, o que exclui $\theta$ = 0. Essa diferença não afeta a nota. Conferi a equivalência algébrica e o exemplo numérico com Python (v $\approx$ 4,38$\times$10$^{5}$ m/s, T $\approx$ 6,56$\times$10$^{-}$$^{7}$ s, cos$\theta$ $\approx$ 0,035, $\theta$ $\approx$ 88$^\circ$). O enunciado transcrito confere com o original (página impressa 9, página 8 do PDF). \\
9 & 10,0 & A resolução da questão 9 está correta e concorda com o gabarito do Poliedro nos três itens (a: equivalente, b: igual, c: equivalente). As formas do campo elétrico coincidem para todo n inteiro (2E0 para n par, 0 para n ímpar), o que conferi numericamente. A interferência é construtiva para n par e destrutiva total para n ímpar. O campo magnético B = -B0 cos(kx - $\omega$t) $\hat y$, com E $\times$ B em +x?, confere pela lei de Faraday. Não há erro em nenhum dos dois. Os pontos levantados são de redação: o caso n = 0 contra d > 0 e o caso n ímpar no item c. \\
10 & 10,0 & A resolução está correta nos dois itens e concorda com o gabarito do Poliedro. Item a: a fórmula de Compton é a mesma, só que na forma $\nu$' = m\_e c$^2$ $\nu$/{[}m\_e c$^2$ + h$\nu$(1 - cos$\theta$){]}. Item b: o observador deve se mover na reta do fóton espalhado, indo ao encontro dele, com V = c($\nu$$^2$ - $\nu$'$^2$)/($\nu$$^2$ + $\nu$'$^2$). Essa expressão é algebricamente igual à do gabarito; conferi as contas numericamente e pelo Doppler relativístico. Também confirmei o enunciado no PDF original (página 10 do arquivo, numerada 11 na prova). As únicas pendências são erros de digitação na resolução detalhada do gabarito: faltam parênteses na expressão final de V e (h$\nu$/c)$^2$ aparece como h$\nu$/c$^2$. Nenhum dos dois afeta a resposta resumida. Resultado: 2 de 2 itens equivalentes ao gabarito. \\
\bottomrule\end{longtable}%
\section*{Principais pontos de dúvida}%
\begin{itemize}%
\item%
\textbf{Q3} \textit{(solution)} Item b: a resolução inclui o impulso de atrito de Coulomb durante o pouso, que vale $\mu$c$\cdot$m$\cdot$v\_z com v\_z = $\surd$(2gh). Com isso o tempo de queda some do resultado: t* = t1 + $\omega$1R/($\mu$c g). O gabarito Poliedro não inclui esse impulso e soma t\_q = $\surd$(2h/g). Os dois modelos são consistentes, e o da resolução é até mais realista para um bloco que cai sobre solo áspero. Mesmo assim, o enunciado ('não quica... mas realiza um deslizamento', sem condição) sugere o modelo simplificado. Cabe ao professor decidir se a resposta da resolução seria aceita, ao menos em parte. A resolução cita t' = t* + $\surd$(2h/g) como 'leitura alternativa', mas não a adota como resposta final.%
\item%
\textbf{Q3} \textit{(statement)} O enunciado não diz o que acontece com a velocidade horizontal no pouso: se ela se conserva ou se o atrito age durante o choque. Também não fornece o coeficiente de atrito cinético entre a massa e o carrossel. Essa ambiguidade é a origem da divergência no item b. Não há gabarito oficial do ITA para confirmar a leitura pretendida.%
\item%
\textbf{Q4} \textit{(statement)} No item b, o enunciado não diz como o gás evolui depois da retirada instantânea da areia. O próprio Poliedro avisa que faltam dados para p\_f em h\_max e dá duas interpretações: adiabática reversível, que leva a uma equação implícita, ou repouso final h\_eq = (1 + 3M/(5m))h\_i. A resolução adota a primeira e a justifica (pistão muito mais lento que o som no gás). Essa é a leitura mais coerente com 'altura máxima' e com 'equação que relaciona'. Mas, se a correção oficial do ITA esperar outra modelagem (por exemplo, expansão irreversível com equilíbrio final em P = mg/A), a nota do item pode mudar. Vale levar a um professor.%
\end{itemize}%
\section*{Questão 1}%
\href{https://poliedroresolve.sistemapoliedro.com.br/vestibulares/ita/2027/ita-2-fase-dia-3-22-10-2026-2-fase-dia-3/dissertativa-1-fisica-2-fase-dia-3-ita-2027}{Resolução do gabarito}\par\smallskip%
\begin{longtable}{@{}p{0.06\textwidth} p{0.33\textwidth} p{0.33\textwidth} p{0.1\textwidth} p{0.08\textwidth}@{}}\toprule Item & Resolução & Gabarito & Veredito & Certo \\ \midrule
a & k\_C = MgL/2 (U''(0) = k - MgL/2; $\omega$$^2$ = (3g/2L)(k/k\_C - 1)) & k\_c = MgL/2 & igual & ambos \\
b & $\lambda$=k/k\_C. Para $\lambda$>1 e para $\lambda$=1: só $\theta$=0, estável (em $\lambda$=1, degenerado, com U-U(0)$\approx$MgL$\theta$$^{4}$/48). Para 2/$\pi$<$\lambda$<1: três equilíbrios, $\theta$=0 instável e $\pm$$\theta$\_e estáveis, com sen$\theta$\_e/$\theta$\_e=$\lambda$. Para 0<$\lambda$$\leq$2/$\pi$: só $\theta$=0 no intervalo aberto, instável. & Para k/k\_c $\geq$ 1: um único equilíbrio, $\theta$=0, estável. Para 2/$\pi$<k/k\_c<1: três equilíbrios, $\theta$=0 instável e $\pm$$\theta$0 estáveis, com k$\theta$0 = k\_c sen$\theta$0. Para k/k\_c $\leq$ 2/$\pi$: um único equilíbrio, $\theta$=0, instável. & igual & ambos \\
c & U = Mg$\cdot$h\_O + (MgL/2){[}cos$\theta$ + $\lambda$$\theta$$^2$/2{]}. Para $\lambda$>1: poço único com mínimo em 0. Para $\lambda$=1: poço quártico com mínimo estrito em 0. Para 2/$\pi$<$\lambda$<1: duplo poço, com máximo local em 0, mínimos em $\pm$$\theta$\_e e barreira MgL$\cdot$sen x(sen x - x cos x), x=$\theta$\_e/2. Para 0<$\lambda$$\leq$2/$\pi$: U decresce de $\theta$=0 até as bordas, sem mínimo no intervalo. Traz quatro esboços ($\lambda$ = 2; 1; 0,8; 0,4). & U($\theta$) = (MgL/2)cos$\theta$ + k$\theta$$^2$/2. Para k $\geq$ k\_c: poço único com mínimo em 0. Para 2/$\pi$<k/k\_c<1: duplo poço, com máximo local em 0 e mínimos em $\pm$$\theta$0. Para k/k\_c $\leq$ 2/$\pi$: máximo em 0 e U decrescendo para as bordas. & equivalente & ambos \\
\bottomrule\end{longtable}%
\begin{itemize}%
\item%
\textbf{a)} Os dois partem do torque resultante (MgL/2)sen$\theta$ - k$\theta$ $\approx$ (MgL/2 - k)$\theta$ e chegam a k\_C = MgL/2. A resolução acrescenta I\_O = ML$^2$/3 e $\omega$$^2$ = (3g/2L)($\lambda$-1). Conferi as duas expressões, e o caso k = 0 dá a taxa de crescimento conhecida, $\surd$(3g/2L). Ela também mostra que em k = k\_C o teste linear não decide.%
\item%
\textbf{b)} As faixas, a quantidade e a natureza dos equilíbrios são as mesmas. A condição sen$\theta$\_e/$\theta$\_e = $\lambda$ é igual a k$\theta$0 = k\_c sen$\theta$0. A resolução é mais rigorosa: prova que f = sen$\theta$/$\theta$ é monótona porque g' = -$\theta$ sen$\theta$ < 0, mostra que U''($\pm$$\theta$\_e) = -k\_C g($\theta$\_e)/$\theta$\_e > 0 e trata à parte o caso degenerado $\lambda$=1. Esse caso é estável porque dU/d$\theta$ = k\_C($\theta$ - sen$\theta$) tem o sinal de $\theta$. Também explica que em $\lambda$ = 2/$\pi$ o equilíbrio cai exatamente em $\theta$ = $\pi$/2, que o intervalo aberto exclui. Conferi numericamente: para $\lambda$ = 0,8, $\theta$\_e = 1,1311 rad (64,8$^\circ$).%
\item%
\textbf{c)} As expressões de U diferem só pela constante Mg$\cdot$h\_O, que desloca o gráfico na vertical e não muda a forma, os extremos nem a estabilidade. As formas dos gráficos coincidem caso a caso. Conferi numericamente os detalhes extras da resolução, e todos estão certos: a barreira para $\lambda$ = 0,8 é 0,0313 MgL; para $\lambda$$\to$1, a barreira tende a (3/4)MgL(1-$\lambda$)$^2$; U($\pi$/2) - U(0) = (MgL/2)($\lambda$$\pi$$^2$/8 - 1) muda de sinal em $\lambda$ = 8/$\pi$$^2$ $\approx$ 0,81; o mínimo plotado vale 0,9374. Os limites de eixo dos quatro esboços comportam as curvas.%
\end{itemize}%
\paragraph{Pontos de dúvida}%
\begin{itemize}%
\item%
\textit{(solution)} O enunciado pede o 'mancal inferior' como referência zero de altura. A resolução põe o eixo O a uma altura constante h\_O acima do mancal e soma Mg$\cdot$h\_O a U. O gabarito põe o eixo na altura zero. A diferença é só uma constante e não muda o esboço, mas a leitura do gabarito (h\_O = 0) parece ser a pretendida.%
\item%
\textit{(solution)} No parágrafo 'Evolução física', a resolução diz que abaixo de (2/$\pi$)k\_C 'não há equilíbrio no intervalo'. É impreciso, porque $\theta$ = 0 continua sendo equilíbrio, só que instável, como a própria resolução afirma no item b. O certo seria 'não há equilíbrio estável no intervalo'. O trecho é complementar e não afeta as respostas.%
\item%
\textit{(key)} Na transcrição do item a do gabarito aparece '$\tau$ $\approx$ MgL/2 - k$\theta$', quando o certo é $\tau$ $\approx$ (MgL/2 - k)$\theta$. Parece erro de digitação ou de extração do texto e não muda o resultado k\_c = MgL/2.%
\item%
\textit{(statement)} O valor de fronteira k/k\_C = 1 não fica claro no enunciado: 'deixa de ser estável para k < k\_C' sugere que em k = k\_C ainda é estável. O gabarito e a resolução concordam que é estável. A resolução justifica isso com o termo de quarta ordem, porque o teste linear não decide. Em k/k\_C = 2/$\pi$, o equilíbrio estável cai em $\theta$ = $\pi$/2, que o intervalo aberto exclui; os dois tratam isso igual.%
\end{itemize}%
\section*{Questão 2}%
\href{https://poliedroresolve.sistemapoliedro.com.br/vestibulares/ita/2027/ita-2-fase-dia-3-22-10-2026-2-fase-dia-3/dissertativa-2-fisica-2-fase-dia-3-ita-2027}{Resolução do gabarito}\par\smallskip%
\begin{longtable}{@{}p{0.06\textwidth} p{0.33\textwidth} p{0.33\textwidth} p{0.1\textwidth} p{0.08\textwidth}@{}}\toprule Item & Resolução & Gabarito & Veredito & Certo \\ \midrule
a & H\_n = e\^{}\{2n\} V$^2$ sen$^2$30$^\circ$/(2g) = (0,49)\^{}n V$^2$/(8g), com e = 0,7 (H\_0 = V$^2$/(8g) no lançamento) & h\_n = (e\^{}n V sen30$^\circ$)$^2$/(2g) = 0,7\^{}(2n) V$^2$/(8g) & igual & ambos \\
b & $\Delta$t = V/{[}g(1-0,7){]} = 10V/(3g) $\approx$ 3,33 V/g (energia mínima E\_min = 3mV$^2$/8 no limite de infinitos choques) & T = (2V sen30$^\circ$/g)$\cdot$1/(1-e) = 10V/(3g) & igual & ambos \\
\bottomrule\end{longtable}%
\begin{itemize}%
\item%
\textbf{a)} Mesma expressão: a velocidade vertical após n choques é e\^{}n$\cdot$V$\cdot$sen30$^\circ$ = 0,7\^{}n$\cdot$V/2, e por Torricelli h = u$^2$/(2g) = 0,7\^{}(2n)V$^2$/(8g). Escrever (0,49)\^{}n em vez de 0,7\^{}(2n) é só outra notação.%
\item%
\textbf{b)} Os dois somam os tempos de voo, que formam a série geométrica (2V sen30$^\circ$/g)(1 + e + e$^2$ + \ldots{}) = (V/g)/(1-0,7) = 10V/(3g). A interpretação também coincide: a energia mínima é a horizontal, que não se perde porque o solo é liso, e é atingida quando a velocidade vertical se anula, depois de infinitos choques em tempo finito. A resolução ainda calcula E\_min = (3/8)mV$^2$, o que o gabarito não pede.%
\end{itemize}%
\paragraph{Pontos de dúvida}%
\begin{itemize}%
\item%
\textit{(statement)} O enunciado diz "mínimo da sua energia" sem dizer qual energia. Os dois leram como energia mecânica total. Essa leitura é a única que dá um instante bem definido, já que a potencial se anula em cada choque e a cinética fica mínima em cada ápice. Ainda assim, é uma interpretação.%
\item%
\textit{(statement)} Não fica claro se "altura máxima após n choques" conta o lançamento como n = 0. A expressão é a mesma nos dois documentos, então isso não muda o resultado.%
\item%
\textit{(solution)} A resolução trata a massa como ponto material e supõe o mesmo e = 0,7 em todos os choques, inclusive nos de velocidade muito pequena. Supõe também que, depois do ponto de acumulação t$\infty$, a massa desliza com v\_x constante. Essas hipóteses são padrão e foram declaradas; o gabarito usa as mesmas sem dizer.%
\item%
\textit{(key)} Na resolução do gabarito há erros de digitação na notação: (e\^{}nVsen$\theta$\^{}2) deveria ser (e\^{}n V sen$\theta$)$^2$, e falta pontuação em "0,7\^{}2nV\^{}20,5\^{}2)/2gh". O resultado final 0,7\^{}(2n)V$^2$/(8g) está correto.%
\end{itemize}%
\section*{Questão 3}%
\href{https://poliedroresolve.sistemapoliedro.com.br/vestibulares/ita/2027/ita-2-fase-dia-3-22-10-2026-2-fase-dia-3/dissertativa-3-fisica-2-fase-dia-3-ita-2027}{Resolução do gabarito}\par\smallskip%
\begin{longtable}{@{}p{0.06\textwidth} p{0.33\textwidth} p{0.33\textwidth} p{0.1\textwidth} p{0.08\textwidth}@{}}\toprule Item & Resolução & Gabarito & Veredito & Certo \\ \midrule
a & d = $\surd$(R$^2$ + (2hR/g)$\cdot$$\surd$($\mu$e$^2$g$^2$ - $\alpha$$^2$R$^2$)), para $\mu$e g > $\alpha$R; se $\mu$e g $\leq$ $\alpha$R, d = R & D = R$\cdot$$\surd$(1 + (2h/g)$\cdot$$\surd$($\mu$e$^2$g$^2$/R$^2$ - $\alpha$$^2$)) & equivalente & ambos \\
b & t* = t1 + $\omega$1R/($\mu$c g) = ($\mu$e$^2$g$^2$ - $\alpha$$^2$R$^2$)\^{}(1/4)($\mu$c g + $\alpha$R)/($\alpha$ $\mu$c g $\surd$R), com o impulso de atrito de Coulomb ($\mu$c$\cdot$m$\cdot$v\_z) atuando durante o choque com o solo. Exige $\omega$1R > $\mu$c$\surd$(2gh); senão t* = t1 + $\surd$(2h/g) & t\_p = $\omega$\_e/$\alpha$ + $\surd$(2h/g) + $\omega$\_e R/($\mu$c g), com $\omega$\_e = ($\mu$e$^2$g$^2$/R$^2$ - $\alpha$$^2$)\^{}(1/4) & diferente & gabarito \\
\bottomrule\end{longtable}%
\begin{itemize}%
\item%
\textbf{a)} R$^2$$\cdot${[}1 + (2h/g)$\cdot$$\surd$($\mu$e$^2$g$^2$/R$^2$ - $\alpha$$^2$){]} = R$^2$ + (2hR/g)$\cdot$$\surd$($\mu$e$^2$g$^2$ - $\alpha$$^2$R$^2$). As duas formas são a mesma expressão, e a conta numérica bate (por exemplo 2,6343 m nos dois, com g=10, R=2, h=1,5, $\alpha$=0,5, $\mu$e=0,5). Os dois usam o mesmo modelo: o atrito estático cede quando R$\surd$($\alpha$$^2$+$\omega$$^{4}$) = $\mu$e g, a massa sai pela tangente com v = $\omega$\_e R, cai em t\_q = $\surd$(2h/g), e por Pitágoras D$^2$ = R$^2$ + (v t\_q)$^2$. A resolução acrescenta o caso $\mu$e g $\leq$ $\alpha$R, que o gabarito não trata; isso é um extra e não muda a resposta principal.%
\item%
\textbf{b)} Pela conta, t1 = $\omega$\_e/$\alpha$ e $\omega$1R = $\omega$\_e R, então a diferença entre as respostas é exatamente $\surd$(2h/g). Isso foi conferido numericamente para vários conjuntos de parâmetros (por exemplo 4,695 s contra 5,243 s). Nenhum dos dois erra a conta: eles usam modelos diferentes para o pouso. O gabarito supõe que, no pouso, a massa perde só a velocidade vertical e mantém a horizontal v\_e. Depois ela freia com $\mu$c g, e t\_f = v\_e/($\mu$c g). A resolução aplica f = $\mu$c N também durante o choque: o impulso de atrito $\mu$c$\cdot$m$\cdot$v\_z tira $\mu$c$\surd$(2gh) da velocidade horizontal. Por isso o tempo de queda se cancela no total. Esse modelo (impacto com atrito de Coulomb) é fisicamente defensável. Mesmo assim, a leitura pretendida pelo enunciado é a do gabarito. O enunciado afirma sem condição que a massa 'não quica ao chegar ao solo, mas realiza um deslizamento'. No modelo da resolução isso nem sempre vale: se $\omega$1R $\leq$ $\mu$c$\surd$(2gh), a massa para no próprio choque, e a resolução precisa de um caso extra que contradiz o enunciado. A própria resolução calcula a 'leitura alternativa' t' = t* + $\surd$(2h/g), que coincide com o gabarito, mas escolhe como resposta final a outra. Pelo critério esperado da banca, a resposta final da resolução provavelmente não levaria a nota cheia.%
\end{itemize}%
\paragraph{Pontos de dúvida}%
\begin{itemize}%
\item%
\textit{(solution)} Item b: a resolução inclui o impulso de atrito de Coulomb durante o pouso, que vale $\mu$c$\cdot$m$\cdot$v\_z com v\_z = $\surd$(2gh). Com isso o tempo de queda some do resultado: t* = t1 + $\omega$1R/($\mu$c g). O gabarito Poliedro não inclui esse impulso e soma t\_q = $\surd$(2h/g). Os dois modelos são consistentes, e o da resolução é até mais realista para um bloco que cai sobre solo áspero. Mesmo assim, o enunciado ('não quica... mas realiza um deslizamento', sem condição) sugere o modelo simplificado. Cabe ao professor decidir se a resposta da resolução seria aceita, ao menos em parte. A resolução cita t' = t* + $\surd$(2h/g) como 'leitura alternativa', mas não a adota como resposta final.%
\item%
\textit{(statement)} O enunciado não diz o que acontece com a velocidade horizontal no pouso: se ela se conserva ou se o atrito age durante o choque. Também não fornece o coeficiente de atrito cinético entre a massa e o carrossel. Essa ambiguidade é a origem da divergência no item b. Não há gabarito oficial do ITA para confirmar a leitura pretendida.%
\item%
\textit{(key)} Há typos na resolução do gabarito. Em (I) aparece $\omega$\_e$^2$ = $\surd$($\mu$\_e g/R - $\alpha$$^2$), e o certo é $\surd$($\mu$e$^2$g$^2$/R$^2$ - $\alpha$$^2$). Nas fórmulas finais de D e de t\_P aparece $\mu$\_e g/R$^2$ dentro da raiz, e o certo é $\mu$e$^2$g$^2$/R$^2$. Aparece também 't\_e = $\omega$\_e/a'. Além disso, a letra $\alpha$ é reaproveitada para a desaceleração linear $\mu$c g, o que se confunde com a aceleração angular. As respostas resumidas do gabarito estão corretas.%
\item%
\textit{(solution)} A resolução supõe $\mu$k < $\mu$e para o atrito cinético entre a massa e o carrossel, que não é dado, e afirma que, com $\mu$k $\geq$ $\mu$e, a massa não deixaria a borda. Essa afirmação é discutível: para $\mu$k = $\mu$e, |a\_D| continua crescendo depois de t1. Isso não afeta as respostas. Os casos extras ($\mu$e g $\leq$ $\alpha$R, que dá d = R e t* = $\surd$(2h/g); e $\omega$1R $\leq$ $\mu$c$\surd$(2gh)) vão além do pedido e não constam do gabarito.%
\end{itemize}%
\section*{Questão 4}%
\href{https://poliedroresolve.sistemapoliedro.com.br/vestibulares/ita/2027/ita-2-fase-dia-3-22-10-2026-2-fase-dia-3/dissertativa-4-fisica-2-fase-dia-3-ita-2027}{Resolução do gabarito}\par\smallskip%
\begin{longtable}{@{}p{0.06\textwidth} p{0.33\textwidth} p{0.33\textwidth} p{0.1\textwidth} p{0.08\textwidth}@{}}\toprule Item & Resolução & Gabarito & Veredito & Certo \\ \midrule
a & h\_f = h\_i$\cdot$((m+M)/m)\^{}\{3/5\}, por adiabática reversível (PV\^{}\{5/3\} = cte) entre P\_i = (m+M)g/A e P\_f = mg/A & h\_f = (1 + M/m)\^{}\{3/5\}$\cdot$h\_i (adiabática reversível, $\gamma$ = 5/3) & igual & ambos \\
b & m(h\_max - h\_i) = (3/2)(m+M)$\cdot$h\_i$\cdot${[}1 - (h\_i/h\_max)\^{}\{2/3\}{]}, supondo evolução adiabática reversível (PV\^{}\{5/3\} = cte) durante o movimento do pistão; raiz não trivial única, com h\_max > h\_f & Forma geral: h\_max = (3M+5m)g h\_i/(3p\_f A + 2mg). O Poliedro avisa que p\_f não é dado. Interpretação 1 (adiabática reversível): 3(M+m)(h\_i/h\_max)\^{}\{5/3\}$\cdot$h\_max + 2m$\cdot$h\_max = (3M+5m)h\_i. Interpretação 2 (repouso final após as oscilações): h\_eq = (1 + 3M/(5m))h\_i & equivalente & ambos \\
\bottomrule\end{longtable}%
\begin{itemize}%
\item%
\textbf{a)} (m+M)/m = 1 + M/m, então as expressões são idênticas. A verificação (iv) da resolução também está certa: W = (3/2){[}(m+M)g h\_i - m g h\_f{]} = -$\Delta$U. O enunciado (PDF, página impressa 5) confere com a transcrição.%
\item%
\textbf{b)} Expandindo a equação da resolução: 2m$\cdot$h\_max - 2m$\cdot$h\_i = 3(m+M)h\_i - 3(m+M)h\_i\^{}\{5/3\}h\_max\^{}\{-2/3\}, ou seja, 2m$\cdot$h\_max + 3(m+M)h\_i\^{}\{5/3\}h\_max\^{}\{-2/3\} = (3M+5m)h\_i. Isso é exatamente a interpretação 1 do gabarito, porque (h\_i/h\_max)\^{}\{5/3\}$\cdot$h\_max = h\_i\^{}\{5/3\}h\_max\^{}\{-2/3\}. Conferência numérica (m = 2, M = 3, h\_i = 0,5): a raiz da resolução dá h\_max $\approx$ 1,455 m e anula o resíduo da equação do gabarito. Também conferi: h\_f $\approx$ 0,866 m; a linearização h\_max - h\_i $\approx$ 6M h\_i/(5(m+M)) (com M = 0,01: 0,00300 contra 0,00299); a forma em y = (h\_max/h\_i)\^{}\{1/3\}; e o estudo de F'(x), cujo mínimo fica em x = h\_f/h\_i. A forma geral do gabarito (com p\_f) também está correta: vem de $\Delta$U = -mg(h\_max - h\_i). A interpretação 2 (h\_eq = 0,95 m no exemplo) também tem a conta certa, mas calcula outra grandeza (o repouso final depois de dissipar a energia), não a altura máxima pedida. Como o enunciado pede 'a equação que relaciona h\_max com m, M e h\_i', a leitura mais coerente é a 1, que é a da resolução.%
\end{itemize}%
\paragraph{Pontos de dúvida}%
\begin{itemize}%
\item%
\textit{(statement)} No item b, o enunciado não diz como o gás evolui depois da retirada instantânea da areia. O próprio Poliedro avisa que faltam dados para p\_f em h\_max e dá duas interpretações: adiabática reversível, que leva a uma equação implícita, ou repouso final h\_eq = (1 + 3M/(5m))h\_i. A resolução adota a primeira e a justifica (pistão muito mais lento que o som no gás). Essa é a leitura mais coerente com 'altura máxima' e com 'equação que relaciona'. Mas, se a correção oficial do ITA esperar outra modelagem (por exemplo, expansão irreversível com equilíbrio final em P = mg/A), a nota do item pode mudar. Vale levar a um professor.%
\item%
\textit{(solution)} A hipótese de que a velocidade do pistão é da ordem de m/s, muito menor que a velocidade do som, é heurística: o enunciado não dá valores numéricos. A resolução a declara de forma explícita, o que é bom, mas ela não decorre do enunciado. Também é uma idealização dizer que o pistão oscila sem dissipação entre h\_i e h\_max.%
\item%
\textit{(key)} A resolução do Poliedro, como foi extraída, tem vários erros de transcrição que não mudam a resposta: p\_i = (M + mg)/A no lugar de (M+m)g/A; 'h\_f = 1 + M/m\^{}(3/5)$\cdot$h\_i' sem parênteses; o sinal em '$\Delta$U = 0 - -W\_ext'; e no item b usa h\_f para a posição de altura máxima. As equações finais estão certas.%
\end{itemize}%
\section*{Questão 5}%
\href{https://poliedroresolve.sistemapoliedro.com.br/vestibulares/ita/2027/ita-2-fase-dia-3-22-10-2026-2-fase-dia-3/dissertativa-5-fisica-2-fase-dia-3-ita-2027}{Resolução do gabarito}\par\smallskip%
\begin{longtable}{@{}p{0.06\textwidth} p{0.33\textwidth} p{0.33\textwidth} p{0.1\textwidth} p{0.08\textwidth}@{}}\toprule Item & Resolução & Gabarito & Veredito & Certo \\ \midrule
a & p' = -n2 R p/(n1R + (n2-n1)p) (imagem virtual, do lado de n1); y'$\pm$ = $\pm$ n1 R h/(n1R + (n2-n1)p) (imagem direta, m = n1R/(n1R+(n2-n1)p)) & p' = -n2 p R/(n1 R + (n2 - n1) p) (imagens virtuais no meio n1); y' = $\pm$ n1 R h/(n1 R + (n2 - n1) p) & igual & ambos \\
b & y1 = $\lambda$0(n1D + n2p)/(2 n1 n2 h) $\approx$ $\lambda$0 D/(2 n2 h), com $\lambda$0 no vácuo (ou y1 = $\lambda$1(n1D+n2p)/(2n2h) com $\lambda$1 = $\lambda$0/n1) & y = $\lambda$D/(2 n2 h), com $\lambda$ no vácuo & equivalente & ambos \\
\bottomrule\end{longtable}%
\begin{itemize}%
\item%
\textbf{a)} Mesma convenção (R\_dioptro = -R, já que C está no meio n1, conforme a figura da p. 6 impressa) e mesmas expressões. Conferi: n1/p + n2/p' = -(n2-n1)/R $\Rightarrow$ p' = -n2Rp/(n1R+(n2-n1)p); m = -n1p'/(n2p) = n1R/(n1R+(n2-n1)p). Fiz um traçado de raio exato com n1=1, n2=1,5, R=10, p=5, h=0,01 e obtive a imagem em (-6,00; 0,0080), igual a p' = -6 e m$\cdot$h = 0,008. A resolução chega a y' pelo raio que passa por C, e o gabarito pelo aumento transversal. O resultado é o mesmo.%
\item%
\textbf{b)} A resolução dá a expressão completa, com fontes virtuais a L = D + n2p/n1 da tela e $\lambda$2 = $\lambda$0/n2. Pelo enunciado, p ? D, então essa expressão se reduz exatamente à do gabarito, e a própria resolução escreve essa aproximação ($\approx$ $\lambda$0D/(2n2h)). Resolvi numericamente o caminho óptico exato pelo princípio de Fermat com $\lambda$0=500 nm, n1=1, n2=1,5, p=1 cm, D=1 m, h=0,1 mm. A 1ª franja fica em 1,69167 mm, que coincide com a fórmula completa da resolução (1,69167 mm). A fórmula do gabarito dá 1,6667 mm, diferença de 1,5\% (ordem n2p/(n1D), desprezível no regime p ? D). As duas adotam $\lambda$ no vácuo.%
\end{itemize}%
\paragraph{Pontos de dúvida}%
\begin{itemize}%
\item%
\textit{(statement)} O enunciado não diz em que meio $\lambda$ é medido. Como as fontes estão em n1, pode{-}se ler $\lambda$ como o comprimento de onda em n1, e então a resposta seria y = n1$\lambda$D/(2n2h). O gabarito supõe $\lambda$ no vácuo e a resolução dá as duas formas, com a do vácuo como principal. As duas concordam, mas a correção oficial pode aceitar ou exigir outra leitura.%
\item%
\textit{(solution)} No item b, a resolução apresenta como resposta principal a fórmula com o termo de correção n2p/n1 (L = D + n2p/n1), mais precisa que a do gabarito, e indica a aproximação p ? D. A diferença é de ordem p/D, que o enunciado manda desprezar. Não deve ser penalizada, mas convém destacar $\lambda$D/(2n2h) como resultado final.%
\item%
\textit{(statement)} O item a pede as posições verticais 'em função de n1, n2, R e p', mas elas dependem também de h, a posição das fontes. Tanto o gabarito quanto a resolução incluem h, o que é inevitável.%
\item%
\textit{(key)} No item a, o sinal de p' depende da convenção (R\_dioptro = -R e p' negativo para imagem virtual). Gabarito e resolução usam a mesma convenção e informam |p'| ou o lado da imagem. Uma resposta com o sinal oposto e a interpretação correta também seria aceitável.%
\end{itemize}%
\section*{Questão 6}%
\href{https://poliedroresolve.sistemapoliedro.com.br/vestibulares/ita/2027/ita-2-fase-dia-3-22-10-2026-2-fase-dia-3/dissertativa-6-fisica-2-fase-dia-3-ita-2027}{Resolução do gabarito}\par\smallskip%
\begin{longtable}{@{}p{0.06\textwidth} p{0.33\textwidth} p{0.33\textwidth} p{0.1\textwidth} p{0.08\textwidth}@{}}\toprule Item & Resolução & Gabarito & Veredito & Certo \\ \midrule
a & q = (3/4)C$\varepsilon$, a carga que sai da armadura superior de C para a placa 1 (ficam C$\varepsilon$/4 em C) & q = 3C$\varepsilon$/4 (C' = 3C, V3 = $\varepsilon$/4) & igual & ambos \\
b & Compressão $\Delta$ = (2/3)d; as placas ficam a d/3 & Compressão $\Delta$x = 2d/3 (distância final d/3) & igual & ambos \\
c & k = 27C$\varepsilon$$^2$/(64d$^2$) = 27$\varepsilon$0a$^2$$\varepsilon$$^2$/(64d$^3$) & k = 27C$\varepsilon$$^2$/(64d$^2$) & igual & ambos \\
\bottomrule\end{longtable}%
\begin{itemize}%
\item%
\textbf{a)} Refiz a conta. Pela conservação de carga no nó isolado, C$\varepsilon$ = (C + C')$\cdot$$\varepsilon$/4, o que dá C' = 3C. Então q = C'$\cdot$V3 = 3C$\varepsilon$/4. A resolução chega ao mesmo valor por outro caminho, q = C($\varepsilon$ - V3). Conferi também com frações exatas.%
\item%
\textbf{b)} Com C' = $\varepsilon$0A/x = 3C = 3$\varepsilon$0A/d, sai x = d/3, e a mola encurta d - d/3 = 2d/3. A resolução usa V(x) = $\varepsilon$x/(x+d) = $\varepsilon$/4 e chega ao mesmo resultado. Os dois dizem que é compressão.%
\item%
\textbf{c)} A força sobre cada placa é F = q$^2$/(2$\varepsilon$0A) = (3C$\varepsilon$/4)$^2$/(2Cd) = 9C$\varepsilon$$^2$/(32d). No equilíbrio, k$\cdot$(2d/3) = F, logo k = 27C$\varepsilon$$^2$/(64d$^2$). Os dois usam o campo da outra placa, $\sigma$/(2$\varepsilon$0), e não 2F. A forma em $\varepsilon$0a$^2$/d$^3$ que a resolução também dá é equivalente, porque $\varepsilon$0A = Cd. A resolução ainda confere a força pela derivada da energia, dU\_e/dx, e o resultado bate.%
\end{itemize}%
\paragraph{Pontos de dúvida}%
\begin{itemize}%
\item%
\textit{(statement)} O valor de k obtido é exatamente o crítico, k\_c = 27C$\varepsilon$$^2$/(64d$^2$). Nele, U'(d/3) = U''(d/3) = 0, ou seja, o equilíbrio em t3 é marginal: é o limiar entre ter e não ter equilíbrio. Com a mola ideal e sem dissipação, as placas soltas em x = d passariam por d/3 e colapsariam. O enunciado só faz sentido se houver amortecimento mecânico. A resolução aponta isso e o gabarito não comenta. As respostas não mudam, mas vale levar a um professor como curiosidade da questão.%
\item%
\textit{(solution)} A resolução diz que parar exatamente em d/3 exige "dissipação mecânica ajustada". A afirmação é imprecisa. Numa dinâmica superamortecida (tipo gradiente), U' $\geq$ 0 com zero único em d/3 faz o sistema se aproximar de d/3 assintoticamente, sem ajuste fino. Basta dissipação forte o suficiente para não ultrapassar esse ponto. Isso é um comentário lateral e não afeta os itens.%
\item%
\textit{(key)} A resolução do gabarito está com a formatação corrompida, provavelmente por perda de frações e quebras de linha na extração. Por exemplo, "C$\varepsilon$ = C + C'$\varepsilon$/4C' = 3C" deveria ser "C$\varepsilon$ = (C + C')$\varepsilon$/4 $\Rightarrow$ C' = 3C", e "k2d/3 = 1/2$\varepsilon$\_0A3C$\varepsilon$/4\^{}2" deveria ser "k$\cdot$(2d/3) = (3C$\varepsilon$/4)$^2$/(2$\varepsilon$0A)". Lidas corretamente, as contas estão certas.%
\item%
\textit{(statement)} A tarefa indica a página 7, que é a página impressa. No arquivo PDF, a Questão 6 está na página 6, e a página 7 do PDF é a Questão 7. Conferi a transcrição contra a página certa e ela é fiel ao original.%
\end{itemize}%
\section*{Questão 7}%
\href{https://poliedroresolve.sistemapoliedro.com.br/vestibulares/ita/2027/ita-2-fase-dia-3-22-10-2026-2-fase-dia-3/dissertativa-7-fisica-2-fase-dia-3-ita-2027}{Resolução do gabarito}\par\smallskip%
\begin{longtable}{@{}p{0.06\textwidth} p{0.33\textwidth} p{0.33\textwidth} p{0.1\textwidth} p{0.08\textwidth}@{}}\toprule Item & Resolução & Gabarito & Veredito & Certo \\ \midrule
a & |V\_o| = (V\_s/4)|G$\varepsilon$ + $\alpha$$\Delta$T| (sinal - para o extensômetro em R1 ou R4, + em R2 ou R3) & |V0| = (Vs/4)$\cdot$|G$\varepsilon$ + $\alpha$$\Delta$T| & igual & ambos \\
b & Quatro extensômetros e nenhum resistor R0: R1 (AD) e R4 (CB) em compressão; R2 (AB) e R3 (DC) em tração (ou tudo trocado). V\_o = V\_sG$\varepsilon$/(1+$\alpha$$\Delta$T) $\approx$ V\_sG$\varepsilon$, sem efeito térmico. & Quatro extensômetros: R2 e R3 na tração, R1 e R4 na compressão (ou tudo trocado), dando V0 = Vs G $\varepsilon$, sem efeito térmico & igual & ambos \\
\bottomrule\end{longtable}%
\begin{itemize}%
\item%
\textbf{a)} Mesma expressão. Refiz a conta: V\_D = V\_s R3/(R1+R3) e V\_B = V\_s R4/(R2+R4), que linearizados dão V\_o $\approx$ (V\_s/4)(x2+x3-x1-x4). Com um único x\_k = G$\varepsilon$+$\alpha$$\Delta$T, sai |V\_o| = (V\_s/4)|G$\varepsilon$+$\alpha$$\Delta$T|. Conferi numericamente (V\_s=10 V, G=2, $\varepsilon$=1e{-}4, $\alpha$$\Delta$T=3e{-}5): a expressão exata dá 0,57493 mV e a linear dá 0,575 mV, com o sinal de cada braço conforme a resolução indica. O detalhe de sinal por braço vai além do que se pede e está correto.%
\item%
\textbf{b)} A configuração e o resultado são os mesmos. Com x2 = x3 = G$\varepsilon$+$\alpha$$\Delta$T e x1 = x4 = -G$\varepsilon$+$\alpha$$\Delta$T, a forma linear dá V\_o = V\_sG$\varepsilon$. A forma exata V\_sG$\varepsilon$/(1+$\alpha$$\Delta$T), que a resolução apresenta, também confere: como x1+x3 = x2+x4 = 2$\alpha$$\Delta$T, o denominador é comum. Numericamente, a exata dá 1,99994 mV, contra 2 mV da linear, e V\_o = 0 para $\varepsilon$ = 0 com qualquer $\Delta$T. A resolução também examina k = 1, 2, 3 e 4 extensômetros e justifica por que a ponte completa é a ótima, o que o gabarito só afirma.%
\end{itemize}%
\paragraph{Pontos de dúvida}%
\begin{itemize}%
\item%
\textit{(statement)} A polaridade da fonte vem só do símbolo da bateria: o traço longo, em cima, é ligado a A, logo V\_A = V\_s. Resolução e gabarito adotam V\_A = V\_s. Com a polaridade invertida, só o sinal de V\_o muda; isso não afeta (a), que pede o módulo, nem a escolha de (b).%
\item%
\textit{(statement)} O enunciado de (b) pede para indicar, para cada dispositivo, se ele é extensômetro ou resistor R0. Isso pode sugerir uma resposta mista, como meia ponte ou extensômetro de compensação ('dummy'). Com pontos de tração e compressão de mesma intensidade, porém, a ponte completa com 4 extensômetros dá sensibilidade máxima (4$\times$ a do item a) e cancela o efeito térmico. Uma meia ponte com braços adjacentes em tração/compressão também cancela $\Delta$T, mas dá só metade do sinal. A correção deveria aceitar apenas a ponte completa como 'máxima precisão'.%
\item%
\textit{(solution)} A resolução declara uma interpretação de 'máxima precisão': erro térmico nulo com sensibilidade máxima. Também assume extensômetros casados (G, $\alpha$ e R0 iguais) e $\Delta$T uniforme. São hipóteses razoáveis, coerentes com 'todo o sistema está na mesma temperatura', e o gabarito usa as mesmas implicitamente.%
\end{itemize}%
\section*{Questão 8}%
\href{https://poliedroresolve.sistemapoliedro.com.br/vestibulares/ita/2027/ita-2-fase-dia-3-22-10-2026-2-fase-dia-3/dissertativa-8-fisica-2-fase-dia-3-ita-2027}{Resolução do gabarito}\par\smallskip%
\begin{longtable}{@{}p{0.06\textwidth} p{0.33\textwidth} p{0.33\textwidth} p{0.1\textwidth} p{0.08\textwidth}@{}}\toprule Item & Resolução & Gabarito & Veredito & Certo \\ \midrule
a & v = $\surd$(2qV/m) & v = $\surd$(2qV/m) & igual & ambos \\
b & Movimento helicoidal uniforme: MRU com velocidade v cos$\theta$ ao longo de B, somado a MCU de raio mv sen$\theta$/(qB) e período 2$\pi$m/(qB) no plano perpendicular a B & Movimento helicoidal uniforme: MCU no plano perpendicular a B (R = mv sen$\theta$/(qB), T = 2$\pi$m/(qB)) composto com MU ao longo de B (v cos$\theta$) & igual & ambos \\
c & cos$\theta$ = (BL/(2$\pi$n))$\surd$(q/(2mV)), com n = 1, 2, 3, \ldots{}, 0 < cos$\theta$ < 1, ou seja, n > (BL/(2$\pi$))$\surd$(q/(2mV)) & cos$\theta$ = qBL/(2$\pi$nmv) = (BL/(2$\pi$n))$\surd$(q/(2mV)), n $\in$ $\mathbb{N}$*, com cos$\theta$ $\leq$ 1, ou seja, n $\geq$ (BL/(2$\pi$))$\surd$(q/(2mV)) & igual & ambos \\
\bottomrule\end{longtable}%
\begin{itemize}%
\item%
\textbf{a)} A resposta sai do teorema da energia cinética: qV = mv$^2$/2. A resolução acrescenta, corretamente, que o módulo da velocidade fica constante dentro do campo, porque a força magnética é perpendicular à velocidade e não realiza trabalho.%
\item%
\textbf{b)} A conclusão, os parâmetros e a justificativa são os mesmos: não há força na direção de B, e a força perpendicular atua como resultante centrípeta. A resolução resolve as equações de movimento explicitamente, informa o passo da hélice, p = 2$\pi$mv cos$\theta$/(qB), e trata os casos degenerados ($\theta$ = 0, $\pi$/2 e $\pi$).%
\item%
\textbf{c)} A expressão é a mesma: L = n$\cdot$v cos$\theta$$\cdot$T, com T = 2$\pi$m/(qB) e v do item a. A forma qBL/(2$\pi$n$\surd$(2qmV)) equivale a (BL/(2$\pi$n))$\surd$(q/(2mV)), o que conferi numericamente. A resolução também mostra que a partícula só volta à linha do ponto 1 em t = nT: a coordenada transversal é y = r(cos $\omega$t - 1), que só se anula quando $\omega$t = 2$\pi$n. A diferença entre a desigualdade estrita da resolução e a não estrita do gabarito é só de fronteira. O caso cos$\theta$ = 1 corresponde a $\theta$ = 0, em que a partícula nunca sai da linha. A desigualdade estrita é um pouco mais coerente com o enunciado, mas a diferença não altera a resposta.%
\end{itemize}%
\paragraph{Pontos de dúvida}%
\begin{itemize}%
\item%
\textit{(key)} Na observação do item c, o gabarito admite cos$\theta$ $\leq$ 1 (n $\geq$ BL/(2$\pi$)$\cdot$$\surd$(q/(2mV))). O caso de igualdade dá $\theta$ = 0: a partícula anda sobre a própria linha e não chega a 'voltar' a ela, como o enunciado pressupõe. A resolução usa a desigualdade estrita. A diferença é só de fronteira e não deve afetar a nota.%
\item%
\textit{(statement)} O enunciado diz 'em um desses retornos', o que justifica a solução geral com n inteiro positivo, adotada pelos dois. Um corretor que esperasse só o primeiro retorno (n = 1) aceitaria a mesma expressão com n = 1. Não há divergência real.%
\end{itemize}%
\section*{Questão 9}%
\href{https://poliedroresolve.sistemapoliedro.com.br/vestibulares/ita/2027/ita-2-fase-dia-3-22-10-2026-2-fase-dia-3/dissertativa-9-fisica-2-fase-dia-3-ita-2027}{Resolução do gabarito}\par\smallskip%
\begin{longtable}{@{}p{0.06\textwidth} p{0.33\textwidth} p{0.33\textwidth} p{0.1\textwidth} p{0.08\textwidth}@{}}\toprule Item & Resolução & Gabarito & Veredito & Certo \\ \midrule
a & E = 2E0 cos(n$\pi$/2) cos(kx - $\omega$t + n$\pi$/2) $\hat z$: onda progressiva no sentido +x, com amplitude 2E0|cos(n$\pi$/2)|, que vale 2E0 para n par e 0 para n ímpar. Para n par, a resolução também simplifica para E = 2E0 cos(kx - $\omega$t) $\hat z$. & E = E0{[}1 + (-1)\^{}n{]} cos(kx - $\omega$t) $\hat z$; amplitude E0{[}1 + (-1)\^{}n{]}; onda no sentido +x & equivalente & ambos \\
b & $\Delta$$\varphi$ = kd = n$\pi$. Para n par a interferência é construtiva (amplitude 2E0, d múltiplo inteiro de $\lambda$). Para n ímpar é destrutiva total (amplitude 0, d múltiplo ímpar de $\lambda$/2). & n par: construtiva (amplitude 2E0); n ímpar: destrutiva total (amplitude 0) & igual & ambos \\
c & B = -(2E0 cos(n$\pi$/2)/c) cos(kx - $\omega$t + n$\pi$/2) $\hat y$ = -(E\_z/c) $\hat y$, com B0 = (2E0/c)|cos(n$\pi$/2)|. Para n par, B = -B0 cos(kx - $\omega$t) $\hat y$ com B0 = 2E0/c, orientado em -$\hat y$. Para n ímpar, B = 0. & B = -B0 cos(kx - $\omega$t) $\hat y$ (E $\times$ B no sentido +x?) & equivalente & ambos \\
\bottomrule\end{longtable}%
\begin{itemize}%
\item%
\textbf{a)} As duas formas são idênticas para todo n inteiro. Testei numericamente para n = 0\ldots{}7 em pontos (x, t, k, $\omega$) aleatórios e a diferença ficou em torno de 1e-15. Para n par vale cos(n$\pi$/2)$\cdot$cos($\theta$ + n$\pi$/2) = (-1)\^{}\{n/2\}$\cdot$(-1)\^{}\{n/2\} cos $\theta$ = cos $\theta$, e a amplitude é 2E0 = E0(1+1). Para n ímpar, cos(n$\pi$/2) = 0 e E0(1-1) = 0. O gabarito chega lá por cos($\theta$ + n$\pi$) = (-1)\^{}n cos $\theta$, e a resolução pela identidade soma{-}produto. As duas mostram que o argumento depende de kx - $\omega$t, ou seja, propagação em +x. A amplitude depende só de E0 e n, como o item pede.%
\item%
\textbf{b)} A conclusão é a mesma do gabarito. Conferi a passagem extra da resolução: com k = 2$\pi$/$\lambda$, d = n$\lambda$/2. Para n par isso dá d = (n/2)$\lambda$, múltiplo inteiro de $\lambda$. Para n ímpar dá um múltiplo ímpar de $\lambda$/2. Está correto.%
\item%
\textbf{c)} A orientação e o sinal coincidem: B está em -$\hat y$ e em fase com E\_z, e $\hat z$ $\times$ (-$\hat y$) = +x?. Refiz pela lei de Faraday: (?$\times$E)\_y = -$\partial$E\_z/$\partial$x = -$\partial$B\_y/$\partial$t, logo B\_y = -E\_z/c. A conferência numérica por diferenças finitas deu resíduo de \textasciitilde{}1e-11. A resolução escreve a forma geral em E0 e n e, para n par, chega exatamente à expressão do gabarito em termos de B0. O caso n ímpar (B = 0) é coerente com o item a; o gabarito deixa esse caso implícito.%
\end{itemize}%
\paragraph{Pontos de dúvida}%
\begin{itemize}%
\item%
\textit{(statement)} O enunciado diz que d é uma constante positiva, mas permite n = 0 em kd = n$\pi$, o que daria d = 0. A resolução declara isso, exclui n = 0 (n $\geq$ 1) e trata n = 0 só como caso{-}limite construtivo. O gabarito não comenta. Nenhuma das duas escolhas muda as conclusões.%
\item%
\textit{(key)} No item c, o gabarito escreve B = -B0 cos(kx - $\omega$t) $\hat y$ sem citar que, para n ímpar, B = 0, nem que B0 = E0{[}1+(-1)\^{}n{]}/c. Isso está implícito em considerar B0 como a amplitude dada. A resolução dá o caso geral (B0 = (2E0/c)|cos(n$\pi$/2)|) e o caso n par na forma do gabarito. Pode valer confirmar com o professor se, no item c, a banca espera a resposta só em função de B0, como o enunciado sugere, ou também a relação B0 = E/c.%
\item%
\textit{(solution)} Na forma geral, a resolução deixa a fase n$\pi$/2 dentro do cosseno e o fator cos(n$\pi$/2) fora: E = 2E0 cos(n$\pi$/2) cos(kx - $\omega$t + n$\pi$/2). Está correto e a simplificação para n par vem logo depois, mas a forma do gabarito, E0{[}1+(-1)\^{}n{]} cos(kx - $\omega$t), atende mais diretamente ao pedido de amplitude expressa só em E0 e n. É questão de apresentação, não de erro.%
\end{itemize}%
\section*{Questão 10}%
\href{https://poliedroresolve.sistemapoliedro.com.br/vestibulares/ita/2027/ita-2-fase-dia-3-22-10-2026-2-fase-dia-3/dissertativa-10-fisica-2-fase-dia-3-ita-2027}{Resolução do gabarito}\par\smallskip%
\begin{longtable}{@{}p{0.06\textwidth} p{0.33\textwidth} p{0.33\textwidth} p{0.1\textwidth} p{0.08\textwidth}@{}}\toprule Item & Resolução & Gabarito & Veredito & Certo \\ \midrule
a & $\nu$' = m\_e c$^2$ $\nu$ / {[}m\_e c$^2$ + h$\nu$(1 - cos$\theta$){]}, equivalente a 1/$\nu$' - 1/$\nu$ = (h/m\_e c$^2$)(1 - cos$\theta$) & $\nu$' = $\nu$ / {[}1 + (h$\nu$/(m\_e c$^2$))(1 - cos$\theta$){]} (Compton) & equivalente & ambos \\
b & |V| = c($\nu$$^2$ - $\nu$'$^2$)/($\nu$$^2$ + $\nu$'$^2$) = c$\cdot$y(2m\_e c$^2$ + y)/(2m\_e$^2$c$^{4}$ + 2m\_e c$^2$ y + y$^2$), com y = h$\nu$(1 - cos$\theta$). O observador fica na reta do fóton espalhado e se move em sentido oposto ao dele, indo ao seu encontro. & O observador se move na direção do fóton espalhado, em sentido oposto ao de sua propagação (aproximando{-}se), com $\beta$ = ($\nu$$^2$ - $\nu$'$^2$)/($\nu$$^2$ + $\nu$'$^2$), ou seja, V = c$\cdot${[}(1 + x)$^2$ - 1{]}/{[}(1 + x)$^2$ + 1{]}, com x = h$\nu$(1 - cos$\theta$)/(m\_e c$^2$). & equivalente & ambos \\
\bottomrule\end{longtable}%
\begin{itemize}%
\item%
\textbf{a)} Basta dividir numerador e denominador da expressão da resolução por m\_e c$^2$ para obter a do gabarito. Refiz a dedução: com E\_e = m\_e c$^2$ + h($\nu$ - $\nu$') e p\_e$^2$c$^2$ = h$^2$($\nu$$^2$ + $\nu$'$^2$ - 2$\nu$$\nu$' cos$\theta$), a relação E\_e$^2$ = p\_e$^2$c$^2$ + m\_e$^2$c$^{4}$ dá m\_e c$^2$($\nu$ - $\nu$') = h$\nu$$\nu$'(1 - cos$\theta$). Também conferi numericamente com valores aleatórios: o resíduo E\_e$^2$ - p\_e$^2$c$^2$ - m\_e$^2$c$^{4}$ fica em \textasciitilde{}1e{-}15. A dedução da resolução é igual à do gabarito.%
\item%
\textbf{b)} Com x = y/M e M = m\_e c$^2$, vale (1 + x)$^2$ - 1 = y(2M + y)/M$^2$ e (1 + x)$^2$ + 1 = (2M$^2$ + 2My + y$^2$)/M$^2$, então as duas expressões coincidem. Conferi numericamente: a diferença Vkey - Vsol fica em \textasciitilde{}1e{-}16. Pelo Doppler relativístico com aproximação, $\nu$'$\cdot$$\surd$((1 + $\beta$)/(1 - $\beta$)) devolve $\nu$ (resíduo \textasciitilde{}1e{-}16). Os sentidos também concordam: o observador se aproxima, indo contra a propagação do fóton espalhado. A diferença de sinal é só convenção. A resolução toma $\beta$ > 0 no sentido do fóton, obtém $\beta$ < 0 e informa |V|. O gabarito já define $\beta$ > 0 como aproximação. Os casos{-}teste da resolução também conferem: $\theta$ = 90$^\circ$ com h$\nu$ = m\_e c$^2$ dá V = 0,6c, e 17,4 keV a 90$^\circ$ dá $\beta$ $\approx$ 0,0335, ou V $\approx$ 1,0$\cdot$10$^{7}$ m/s.%
\end{itemize}%
\paragraph{Pontos de dúvida}%
\begin{itemize}%
\item%
\textit{(key)} Na resolução detalhada do gabarito, a expressão final de V aparece como (1 + h$\nu$/(m\_ec$^2$)(1-cos$\theta$)\^{}2 - 1)/(1 + h$\nu$/(m\_ec$^2$)(1-cos$\theta$)\^{}2 + 1)$\cdot$c, sem os parênteses externos. Lida literalmente, ela eleva só (1 - cos$\theta$) ao quadrado e fica errada. A resposta resumida do gabarito, {[}(1 + h$\nu$(1-cos$\theta$)/(m\_e c$^2$))$^2$ - 1{]}/{[}(1 + \ldots{})$^2$ + 1{]}, está correta, então o problema parece ser só de digitação ou extração.%
\item%
\textit{(key)} Também na resolução do gabarito, a linha da lei dos cossenos está escrita como p\_e$^2$ = h$\nu$/c$^2$ + (h$\nu$')/c$^2$ - \ldots{}. O certo é (h$\nu$/c)$^2$ + (h$\nu$'/c)$^2$ - 2(h$\nu$/c)(h$\nu$'/c)cos$\theta$. Parece erro de transcrição e não afeta o resultado.%
\item%
\textit{(statement)} O enunciado diz que o observador se move 'na mesma direção do fóton espalhado', o que deixa o sentido em aberto. Como $\nu$' < $\nu$, só existe solução com o observador indo ao encontro do fóton, em sentido oposto ao de propagação. A resolução e o gabarito interpretam da mesma forma, e a resolução explica isso numa observação.%
\item%
\textit{(solution)} Na convenção da resolução, $\beta$ > 0 é o sentido de propagação do fóton espalhado, então a fórmula sai com $\beta$ = ($\nu$'$^2$ - $\nu$$^2$)/($\nu$'$^2$ + $\nu$$^2$) < 0 e a resposta é dada em módulo, com o sentido indicado por extenso. O gabarito usa $\beta$ > 0 para aproximação. É só convenção de sinal, e o resultado físico é o mesmo.%
\end{itemize}%
\end{document}